\section{Concepts I: Instruction-Level Parallelism}\label{grp:ilp}
\textbf{Dependences decide everything here.} True (RAW): real data flow, \emph{cannot} be removed. Anti (WAR) and output (WAW): name conflicts, removed by \emph{renaming}. Control: relaxed by prediction + speculation. Memory dependences are hard to see (does \texttt{100(R4)} equal \texttt{20(R6)}?). Loop-carried dependence (\texttt{A[i+1]=A[i]+C[i]}) stops iterations running in parallel.
{\small\renewcommand{\arraystretch}{1.3}\rowcolors{1}{white}{white}
\begin{longtable}{L{2.3cm}L{4.7cm}L{2.9cm}L{3.0cm}L{2.8cm}}
\toprule
\rowcolor{hdrc}\textbf{Concept} & \textbf{What it fixes $\cdot$ how it works} & \textbf{Best scenario (use when)} & \textbf{Why it beats the alternatives} & \textbf{Weakness $\cdot$ avoid when}\\
\midrule
\endfirsthead
\toprule
\rowcolor{hdrc}\textbf{Concept} & \textbf{What it fixes $\cdot$ how it works} & \textbf{Best scenario (use when)} & \textbf{Why it beats the alternatives} & \textbf{Weakness $\cdot$ avoid when}\\
\midrule
\endhead
\bottomrule
\endfoot
\textbf{\textcolor{cardc}{I1}}\label{card:I1}\newline\textbf{Loop unrolling + static (compiler) scheduling}\newline{\footnotesize\textit{Static}} & \textbf{Fixes:} Latency stalls and loop overhead in regular loops. \newline \textbf{How:} Replicate the body, rename registers (F0/F4, F6/F8\ldots), group loads, then adds, then stores, so independent work fills latency gaps; branch overhead amortised. & Independent iterations, known fixed latencies, many registers, in-order/VLIW/embedded cores, regular FP loops. & Dynamic scheduling here: same gain with no OoO hardware, power or area. & Register pressure, code growth; needs alias knowledge. \newline \textbf{Avoid:} Unknown aliasing, pointer chasing, variable cache-miss latency, binary must run well on future hardware. \\*
\rowcolor{black!4}\multicolumn{5}{p{\dimexpr\linewidth-2\tabcolsep\relax}}{\textbf{Numbers:} Slide loop: 9 cycles/iter $\rightarrow$ 6 (reordered) $\rightarrow$ 6.75 (unrolled $\times$4, unscheduled) $\rightarrow$ \textbf{3.5} (unrolled + scheduled). \enspace$\cdot$\enspace \textbf{\textcolor{pickc}{Write:}} \textit{The loop has independent iterations and fixed latencies [given], so unrolling $\times$4 with renaming fills the 1--3-cycle gaps: 9 $\rightarrow$ 3.5 cycles/iteration without OoO hardware.}}\\
\midrule
\textbf{\textcolor{cardc}{I2}}\label{card:I2}\newline\textbf{Software pipelining}\newline{\footnotesize\textit{Static}} & \textbf{Fixes:} Long FP latencies in loops. \newline \textbf{How:} Each new loop body holds one instruction from each of several iterations (load of $i{+}2$, add of $i{+}1$, store of $i$). & Long FP latencies, VLIW/in-order cores, where unrolling would blow up code size. & Unrolling: less code growth for the same overlap. & Prologue/epilogue code, register pressure. \\*
\midrule
\textbf{\textcolor{cardc}{I3}}\label{card:I3}\newline\textbf{Scoreboard (CDC 6600)}\newline{\footnotesize\textit{Dynamic}} & \textbf{Fixes:} In-order stalls behind one slow instruction. \newline \textbf{How:} In-order issue, out-of-order execution. Issue (FU free, no WAW) $\rightarrow$ read operands (wait for RAW) $\rightarrow$ execute $\rightarrow$ write (wait for WAR). Central tables. & Some OoO tolerance with minimal logic, few units. & In-order pipeline: independent instructions pass a stalled one (\texttt{DIVD F0;\ ADDD F10,F0;\ SUBD F12}: SUBD need not wait). & No forwarding, window $\le$5, stalls on WAR/WAW and structural hazards. \\*
\rowcolor{black!4}\multicolumn{5}{p{\dimexpr\linewidth-2\tabcolsep\relax}}{\textbf{Numbers:} Worked 6-instruction example: finishes cycle 62 (Tomasulo 57). \enspace$\cdot$\enspace \textbf{\textcolor{pickc}{Write:}} \textit{A scoreboard lets [independent instructions] proceed past the [long divide] with only central tables, accepting WAR/WAW stalls because the code has few name conflicts.}}\\
\midrule
\textbf{\textcolor{cardc}{I4}}\label{card:I4}\newline\textbf{Tomasulo (hardware register renaming)}\newline{\footnotesize\textit{Dynamic}} & \textbf{Fixes:} WAR/WAW stalls, few architectural registers, variable latencies. \newline \textbf{How:} Reservation stations at each unit hold operands or tags of their producers (renaming); results broadcast on the common data bus (CDB) to all waiting stations; loads/stores have buffers. Overlaps loop iterations (``dynamic unrolling'') and hides cache misses. & FP code with long/variable latencies, cache misses, only few registers (360/91 had 4 FP regs), code not tuned for this machine. & Scoreboard: renaming removes WAR/WAW stalls; window $\le$14 vs.\ $\le$5. Static scheduling: adapts to run-time latencies. & Complex; associative CDB is the limit (one result per CDB per cycle, high capacitance); imprecise interrupts. \\*
\rowcolor{black!4}\multicolumn{5}{p{\dimexpr\linewidth-2\tabcolsep\relax}}{\textbf{Needs:} ROB for precise exceptions/speculation \cp{I5}; non-blocking cache \cp{M10}. \enspace$\cdot$\enspace \textbf{Numbers:} IBM 360/91: 6 load, 3 store, 3 add, 2 mult/div buffers/stations; CDB $=$ 64 data + 4 tag bits. \enspace$\cdot$\enspace \textbf{\textcolor{pickc}{Write:}} \textit{Only [4] FP registers and variable miss latency [given] defeat static scheduling; Tomasulo renames through reservation stations so WAR/WAW vanish and iterations overlap; its limit is one result per CDB per cycle.}}\\
\midrule
\textbf{\textcolor{cardc}{I5}}\label{card:I5}\newline\textbf{Speculation: Tomasulo + reorder buffer (ROB)}\newline{\footnotesize\textit{Dynamic}} & \textbf{Fixes:} Control dependences limit the window; need precise exceptions. \newline \textbf{How:} Predict branches, execute past them, keep results in the ROB in program order; \textbf{execute out of order, commit in order}. Mispredicted branch at ROB head flushes younger entries. Stores write memory only at commit; a load waits if an earlier store has a matching address. & High-performance single-thread cores, wide issue, branchy code, precise exceptions required. & Non-speculative Tomasulo: schedules across basic blocks; early recovery; precise interrupts. & Most area/power; many ports; wasted work on mispredicts (Lab 7: 20 cycles of wrong-path work); speculation hurt by cache/TLB misses. \\*
\rowcolor{black!4}\multicolumn{5}{p{\dimexpr\linewidth-2\tabcolsep\relax}}{\textbf{Needs:} Accurate predictor \cp{P10} + BTB \cp{P11}; physical registers, LQ/SQ, checkpoints sized together (Lab 7). \enspace$\cdot$\enspace \textbf{Numbers:} Four steps: issue (free RS \emph{and} ROB entry) $\rightarrow$ execute $\rightarrow$ write result (CDB $\rightarrow$ ROB) $\rightarrow$ commit at ROB head. \enspace$\cdot$\enspace \textbf{\textcolor{pickc}{Write:}} \textit{Branchy single-thread code needs instructions from beyond [k] branches; speculation with a ROB keeps the units busy and still gives precise exceptions.}}\\
\midrule
\textbf{\textcolor{cardc}{I6}}\label{card:I6}\newline\textbf{Dynamically scheduled superscalar (multiple issue)}\newline{\footnotesize\textit{Dynamic}} & \textbf{Fixes:} CPI cannot go below 1 with single issue. \newline \textbf{How:} Fetch/decode/issue several instructions per cycle; hardware checks dependences among them; multiple CDBs/commit. & General-purpose products, legacy binaries, irregular code, unpredictable latencies; width $\approx$4. & VLIW: binary compatible across generations, tolerates variable latency, smaller code. & Issue logic is the bottleneck; ports and power grow faster than width. \\*
\rowcolor{black!4}\multicolumn{5}{p{\dimexpr\linewidth-2\tabcolsep\relax}}{\textbf{Needs:} Predictor + BTB, renaming, more register/memory ports, ROB, multiple CDBs/commit, decode width (Lab 6). \enspace$\cdot$\enspace \textbf{Numbers:} 3--6 issue is the practical limit; Athlon (3-issue) led SPECint/transistor and /area. \enspace$\cdot$\enspace \textbf{\textcolor{pickc}{Write:}} \textit{The product must run [legacy binaries] whose latencies vary with cache misses, so a 4-wide dynamically scheduled core finds ILP at run time; VLIW would need recompilation for every generation.}}\\
\midrule
\textbf{\textcolor{cardc}{I7}}\label{card:I7}\newline\textbf{VLIW / EPIC}\newline{\footnotesize\textit{Static}} & \textbf{Fixes:} Issue-logic cost of finding parallelism in hardware. \newline \textbf{How:} Compiler packs $K$ independent operations per long instruction (e.g.\ 1 int/branch, 2 FP, 2 memory). & Embedded/DSP/fixed product, controlled compiler, regular loops with abundant ILP, no compatibility need. & Superscalar there: simpler decode/issue, no interlocks, lower power. & Code size (NOPs), many register ports (example: 7R/3W int, 5R/3W FP). \newline \textbf{Avoid:} Binary compatibility needed (changing units/latencies breaks code), cache misses (lock-step: one stall stalls the bundle), irregular code. \\*
\rowcolor{black!4}\multicolumn{5}{p{\dimexpr\linewidth-2\tabcolsep\relax}}{\textbf{Numbers:} Itanium 2: least efficient per transistor/area/watt except SPECfp/W. \enspace$\cdot$\enspace \textbf{\textcolor{pickc}{Write:}} \textit{The DSP runs fixed loops compiled once and needs no compatibility, so VLIW packs [K] operations per word with no issue logic; a superscalar's dependence-check hardware would cost power for nothing.}}\\
\midrule
\textbf{\textcolor{cardc}{I8}}\label{card:I8}\newline\textbf{Break dependence chains (multiple accumulators, reassociation, tree)}\newline{\footnotesize\textit{Static}} & \textbf{Fixes:} Low IPC from a loop-carried recurrence (\texttt{sum += \ldots}). \newline \textbf{How:} Use $k$ independent partial sums ($k=$ latency $\times$ pipelined units), combine at the end; tree reduction for many values. & Reductions, dot products, accumulations; wide core with idle units. & Wider issue/bigger window: they cannot overlap a true dependence (Lab 8: 8-wide, 96-entry scheduler gave the same 123 cycles; Lab 10: recurrence 0.27 IPC vs.\ $\approx$3 independent). & FP reassociation changes rounding slightly; more registers. \\*
\rowcolor{black!4}\multicolumn{5}{p{\dimexpr\linewidth-2\tabcolsep\relax}}{\textbf{\textcolor{pickc}{Write:}} \textit{The FMA latency is [4] and there are [2] units, so 8 accumulators keep both busy; the chain, not the width, limited IPC to [1.25].}}\\
\midrule
\textbf{\textcolor{cardc}{I9}}\label{card:I9}\newline\textbf{Wider issue / bigger OoO window (when it fails)}\newline{\footnotesize\textit{---}} & \textbf{Fixes:} Only idle slots caused by too few instructions in flight. \newline \textbf{How:} Needs $\approx$depth$\times$FUs independent instructions (15--20); doubling to 6--12 issue needs 3--4 memory accesses, 2--3 branches, 20+ register accesses and 12--24 fetches per cycle. & Only when independent instructions exist and every other structure (predictor, ports, PRF, decode) keeps up. & Nothing when code is dependence-, memory- or branch-bound: there TLP/DLP \cp{E3} or chain-breaking \cp{I8} win. & Clock suffers; energy-inefficient (switching $\propto$ peak, performance $\propto$ sustained); practical windows 32--126. \newline \textbf{Avoid:} dependence-bound, memory-bound or branch-bound code; power limits. \\*
\rowcolor{black!4}\multicolumn{5}{p{\dimexpr\linewidth-2\tabcolsep\relax}}{\textbf{Numbers:} Lab 7: enlarging one structure moves the stall to the next; Lab 8: add units only of the saturated class. \enspace$\cdot$\enspace \textbf{\textcolor{pickc}{Write:}} \textit{Doubling issue would need 3--4 memory ports and 2--3 branch resolutions per cycle, slowing the clock; the [x] idle slots come from [dependences / misses], so I spend the area on [threads / cores] instead.}}\\
\end{longtable}}

\small
\begin{tabularx}{\textwidth}{L{2.6cm}L{3.6cm}L{4.2cm}X}
\toprule
\tblhead{\textbf{} & \textbf{Scoreboard} & \textbf{Tomasulo} & \textbf{Tomasulo + ROB (speculative)}}\\
\midrule
WAR / WAW & Stall write / stall issue & Removed by renaming & Removed by renaming \\
Results & Through registers, no forwarding & Broadcast on CDB & CDB $\rightarrow$ ROB $\rightarrow$ commit \\
Window & $\le$5 & $\le$14 & ROB size (32--126 practical) \\
Exceptions & Imprecise & Imprecise & \textbf{Precise} (at commit) \\
Branches & Stop at branch & Wait for branch & Speculate past branches \\
Pick for & Minimal OoO, few units & FP, few registers, variable latency & High-performance core, wide issue \\
\bottomrule
\end{tabularx}
\normalsize

\small
\begin{tabularx}{\textwidth}{L{2.8cm}L{6.2cm}X}
\toprule
\tblhead{\textbf{} & \textbf{Superscalar (dynamic) \cp{I6}} & \textbf{VLIW (static) \cp{I7}}}\\
\midrule
Finds parallelism & Hardware at run time (forwarding, renaming) & Compiler: bundles of $K$ independent ops \\
Binary compatibility & Yes, across generations & Weak: new units/latencies break schedules (emulation is slow) \\
Variable latency (misses) & Independent instructions continue & Lock-step: one stall stalls the bundle \\
Hardware / power & Complex issue, dependence checks, renaming & Simple decode/issue, no interlocks \\
Code size & Small & Large (NOPs, unrolling) \\
Examples & Pentium 4, Athlon 64, Power5 & Itanium 2 (IA-64/EPIC) \\
\bottomrule
\end{tabularx}
\normalsize
\textbf{ILP efficiency ranks (1 = best; Itanium 2 / Pentium 4 / Athlon 64 / Power5):} SPECint per transistor and per area 4/2/1/3; SPECint/W 4/3/1/2; SPECfp/W 2/4/3/1. No single leader; VLIW ``not a significant breakthrough in scaling ILP''.
